The physics of the dark photon analysis necessarily creates a challenging environment for vertex tagging.
In the ggF production mode, the observable signature is a photon, \ETmiss, and less than four initial state radiation jets.
This is a low track multiplicitiy environment, and tradition vertex tagging techniques are shown to be not as efficient, as described in section~\ref{sec:vertex}.
Briefly, the vertex that is selected as the primary vertex is the one with the highest sum of track \pt squared, $\sum p_{T}^2$.
The method is robust and performs well in high track multiplicity environments, but is shown to be less efficient in a low track multiplicity environment, such as the dark photon analysis. 
In this section, preliminary studies are shown using a neural network classifier approach to improve the vertex tagging efficiency for a future iteration of the dark photon analysis, potentially a full Run 3 search.

\subsection{Current vertex tagging performance and motivation}

As stated above, the standard ATLAS hard-scatter vertex selection chooses the vertex with the highest $\sum p_{T}^2$ of associated tracks as the primary vertex.
In Run 3 with a $\langle \mu \rangle$ of 60, there are multiple reconstructed vertices per event, each built from at least two tracks, and a transverse momentum of at least 400 MeV.
The final state of the dark photon analysis is a photon + \ETmiss, and as a result, the track activity is low.
This is mainly due to the absence of charged leptons or jets in the final state.
Relying on the standard vertex tagging method can lead to a pile-up vertex, the result of multiple soft scatter interactions, being incorrectly identified as the primary vertex since it may have a higher $\sum p_{T}^2$ than the actual hard-scatter vertex.
This misidentification can lead to a significant inaccuracy in the \ETmiss calculation, as seen in Fig.~\ref{fig:etmiss_misidentification}, and can degrade the overall sensitivity of the search for dark photons.
When the primary vertex is misidentified, pile-up jets can be incorrectly associated with the hard-scatter vertex, leading to an overestimation of the \ETmiss.

\begin{figure}[htbp]
    \centering
    \includegraphics[width=0.6\textwidth]{../images/dark_photon/vertexing/fake_met.png}
    \caption{Diagram of a hard scatter jet and photon and a pile-up vertex with a jet, with the hard scatter vertex correctly identified as the primary vertex (left) and incorrectly identified as the pile-up vertex (right). The misidentification of the primary vertex leads to a significant miscalculation of the \ETmiss.}
    \label{fig:etmiss_misidentification}
\end{figure}

How often the primary vertex is misidentified in the dark photon analysis varies, with up to 30 percent of signal events and 90 percent of $\gamma$ + jets background events incorrectly tagged with preselections applied.
Misidentification presents a gaussian distributed background if the reconstructed z position of the primary vertex is plotted against the true z position of the hard-scatter vertex, as shown in Fig.~\ref{fig:vertex_misidentification}.
It appears that either the correct vertex is identified exactly, or a random pile-up vertex is identified, leading to a gaussian distribution around the diagonal line representing perfect vertex tagging.

\begin{figure}[htbp]
    \centering
    \includegraphics[width=0.45\textwidth]{../images/dark_photon/vertexing/pv_baseline_vs_truth_ggHyyd.pdf}
    \caption{The reconstructed z position of the primary vertex plotted against the true z position of the hard-scatter vertex for signal events. The diagonal line represents perfect vertex tagging, while the gaussian distribution around the diagonal line represents misidentification.}
    \label{fig:vertex_misidentification}
\end{figure}    

With a \ETmiss cut of 100 GeV in the signal region, this leads to a significant amount of background events entering the signal region, seen in Fig.~\ref{fig:etmiss_misidentification}, which can degrade the sensitivity of the analysis.
This motivates the need for improved vertex tagging techniques that can better distinguish between the hard-scatter vertex.

\subsection{Preliminary studies with a neural network classifier}

To address the issue of vertex misidentification for the next iteration of the dark photon analysis, preliminary studies have been conducted using a neural network classifier approach.
The goal of this approach is to improve the vertex tagging efficiency by training a neural network to distinguish between the hard-scatter vertex and pile-up vertices based on various features of the vertices, while also taking into account the geometic information from the photon showers in the calorimeter.
The reconstructed photon showers in the first and second layers of the calorimeter are used in combination with the top three primary vertex candidates ordered in $\sum p_{T}^2$.
A linear fit is then performed using the z and R positions of the photon showers and the z positions of the three primary vertex candidates to extract a $\chi^2$ value as an input to the neural network, also called a ``geometric input" to the neural network.
This fitting method can be seen in the diagrams in Fig.~\ref{fig:vertexing_fit}, showing the 2 layers of the EM calorimeter, the photon showers, and the three primary vertex candidates.


\begin{figure}[htbp]
    \centering
    \includegraphics[width=0.45\textwidth]{../images/dark_photon/vertexing/vertexing_fit.png}
    \includegraphics[width=0.45\textwidth]{../images/dark_photon/vertexing/fit_goodreco6.png}
    \caption{Diagram of the linear fit performed using the z and R positions of the photon showers and the z positions of the three primary vertex candidates (left). Actual line of best fit example for a signal event (right). The star is the true vertex, and the $\chi^2$ value from this fit is used as an input to the neural network classifier.}
    \label{fig:vertexing_fit}
\end{figure}

\subsubsection{Neural network architecture and training}
A supervised binary MLP classifier from scikit-learn is used for this preliminary study, and the features used for training are listed in Table~\ref{tab:nn_features}.
This is a feed-forward network with the task of scoring the three primary vertex candidates, with the highest score being selected as the primary vertex.
It is important to note that all features are candidate specific except for the $\Delta E^{miss}_T$ and photon pointing, which are the same for all three candidates in an event.

\begin{table}[htbp]
    \centering
    \small
    \begin{tabular}{|p{0.34\textwidth}|p{0.58\textwidth}|}
        \hline
        Feature & Description \\
        \hline
        $log(\sum p_{T}^2)$ & Logarithm of the sum of the squared transverse momentum of tracks associated with the vertex \\
        $N_{\text{tracks}}$ & Number of tracks associated with the vertex \\
        $z_{\text{vertex}}$ & z position of the vertex \\
        $\gamma$ pointing $= \frac{z_\gamma - z_{\text{vertex}}}{\sigma_{z_\gamma}}$ & Photon pointing variable \\
        $\Delta E^{miss}_T$ & Difference between \ETmiss calculation with and without pile-up jets \\
        $\chi^2$ & Chi-squared value from the linear fit of photon showers and vertex positions \\
        \hline
    \end{tabular}
    \caption{Features used for training the neural network classifier to improve vertex tagging efficiency.}
    \label{tab:nn_features}

\end{table}

The architecture of the neural network consists of two hidden layers with 32 and 16 neurons respectively, using the ReLU activation function, and an output layer with a single neuron using the sigmoid activation function for binary classification.
Two-class softmax is used as the loss function, and the Adam optimizer is employed for training.
The probability for class "True PV" is used as the score for each vertex candidate, and the candidate with the highest score is selected as the primary vertex.
Training is done on 30\% of the available signal and a $p_T$ slice from 70 GeV to 140 GeV of $\gamma$ + jets background events, with the remaining 70\% used for testing and evaluation, with a maximum number of iterations set to 200.
The background slice of events are weighted.
Two sets of training are performed: one on events without selections, and one with preselections applied as defined in section~\ref{sec:selection}, to evaluate the performance of the classifier in both scenarios.
At most, there is one positive candidate per event, closest to the true vertex within .5 mm, else all candidates are labeled as negative.
As seen in Fig.~\ref{pointing}, the photon pointing variable is a powerful discriminator between the hard-scatter vertex and pile-up vertices, as the pointing for incorrectly identified vertices is much broader than the pointing for correctly identified vertices, which is more peaked around zero.
The $log(\sum p_{T}^2)$ and $\chi^2$ variables also provide significant discrimination power as expected.

\begin{figure}[htbp]
    \centering
    \includegraphics[width=0.45\textwidth]{../images/dark_photon/vertexing/feature_ph_vtx_gamma_jets.pdf}
    \caption{Distribution of the photon pointing variable for correctly identified vertices and incorrectly identified vertices in the background slice. The photon pointing variable shows good separation between correctly identified vertices and incorrectly identified vertices, with correctly identified vertices having a more peaked distribution around zero, while incorrectly identified vertices have a broader distribution.}
    \label{pointing}
\end{figure}


\subsubsection{Preliminary results and performance evaluation}

Predicted probabilities from the neural network classifier are used to select the primary vertex candidate with the highest score, and the performance of the classifier is evaluated on the test set.
These probabilities are plotted for candidates split into correctly identified vertices and incorrectly identified vertices, as shown in Fig.~\ref{probabilities}, for both the training with preselections and without preselections.
For the $\gamma$ + jets background slice, the classifier shows excellent separation for no preselections, with the majority of true PV candidates having predicted probabilities close to 1.
With preselections, the separation is still good, especially compared to the stand reconstruction.
In signal events, there is a decent separation between the true PV candidates and pile-up vertices, but there is a large fraction of candidates centered around a probability of 0.2 without preselections.
As expected, the preselections clean up the signal events, leading to a grouping of true PV candidates around higher predicted probabilities, while the pile-up vertices are grouped around zero.


\begin{figure}[htbp]
    \centering
    \includegraphics[width=0.45\textwidth]{../images/dark_photon/vertexing/prob_distribution_gammajet_frag_70_140_weighted.pdf}
    \includegraphics[width=0.45\textwidth]{../images/dark_photon/vertexing/prob_distribution_ggHyyd.pdf}
    \includegraphics[width=0.45\textwidth]{../images/dark_photon/vertexing/prob_distribution_gammajet_frag_70_140_presel_weighted.pdf}
    \includegraphics[width=0.45\textwidth]{../images/dark_photon/vertexing/prob_distribution_ggHyyd_presel.pdf}
    \caption{Predicted probabilities from the neural network classifier for candidates split into correctly identified vertices and incorrectly identified vertices, for both the training with preselections (bottom) and without preselections (top), and split by background (left) and signal (right). The classifier shows good separation between the true PV candidates and pile-up vertices, especially in the background slice.}
    \label{probabilities}
\end{figure}

The final performance of the classifier is evaluated using a simple accuracy metric, which is the percentage of events where the primary vertex is correctly identified.
These percentages are listed in Table~\ref{tab:nn_performance} for both the training with preselections and without preselections, and compared to the performance of the standard vertex tagging method. % this table is actually a png from another pdf generated from a latex table. 
With preselections applied, the classifer shows a significant improvement in vertex tagging efficiency, especially in the background slice, where the efficiency increases from ~20\% to ~51\%. 
In signal events, the efficiency also improves from ~65\% to ~72\%, which is a significant improvement.
For a preliminary study, these results are very promising, especially in the background rejection potential of the classifier.

\begin{figure}[htbp]
    \centering
    \includegraphics[width=0.9\textwidth]{../images/dark_photon/vertexing/table_results.png}
    \caption{Table of the vertex tagging efficiency for the standard vertex tagging method and the neural network classifier, for both the training with preselections and without preselections, and split by background and signal. The neural network classifier shows a significant improvement in vertex tagging efficiency compared to the standard method, especially in the background slice.}
    \label{tab:nn_performance}
\end{figure}

The $\Delta_z$ between the reconstructed primary vertex and the true vertex is also plotted for both the standard vertex tagging method and the neural network classifier, as shown in Fig.~\ref{delta_z}, for the training with preselections, and split by background and signal.
There is a significant improvement in the $\Delta_z$ distribution for the neural network classifier, with a sharper peak around zero for correctly identified vertices, especially in the background slice, indicating that the classifier is more accurately identifying the primary vertex compared to the standard method.

\begin{figure}[htbp]
    \centering
    \includegraphics[width=0.45\textwidth]{../images/dark_photon/vertexing/delta_z_combined_gammajet_vs_signal.pdf}
    \caption{The $\Delta_z$ between the reconstructed primary vertex and the true vertex for both the standard vertex tagging method (solid) and the neural network classifier (dashed) with preselections applied, and split by background (red) and signal (blue). The neural network classifier shows a significant improvement in the $\Delta_z$ distribution, with a sharper peak around zero for correctly identified vertices, especially in the background slice.}
    \label{delta_z}
\end{figure}

\subsection{Future work and conclusions}

While these preliminary results are promising, there is still work to be done to further optimize the neural network classifier and evaluate its performance in a full analysis context.
Future work includes exploring different neural network architectures, such as a Boosted Decision Tree (BDT) or a more complex deep learning model, and performing a more detailed evaluation of the classifier's performance in terms of its impact on the overall sensitivity of the dark photon analysis by including more backgrounds.
Additionally, a more comprehensive metric for performance evaluation would be to run the full particle flow reconstruction and evaluate the \ETmiss resolution and the impact on the signal significance in the analysis.
With these preliminary studies, a conservative estimate for the improvement in signal sensitivity for the dark photon analysis with the neural network classifier is around ~10\%.
This is based on the asymptotic scaling of the signal significance with the square root of background reduction over the increase in signal events.
In conclusion, these preliminary studies show that a neural network approach to vertex tagging has the potential to significantly improve the vertex tagging efficiency in the dark photon analysis for a full Run 3 search.
